\section{Discussion}
\label{sec:discussion}

\textbf{What the comparison reveals.} The symbolic baseline and the
neuro-symbolic analyzer fail in different ways. PyCG-2026 is \emph{exact}
where the semantics is fully specified and it has a model --- argument binding,
MRO, typeshed-typed outside values, \texttt{match} --- and it degrades where a
value depends on information it does not interpret: string-valued dispatch,
\texttt{eval}, documented roles, library callbacks. \nscg reads code the way a
developer does and handles string-built dispatch and dynamic features well, but
it over-approximates with the hints it is given, misnames outside callees, and,
crucially, it does not \emph{discover} callees that no symbolic component
proposed. The symbolic stages are therefore not just a guard rail: they define
the search space of the neural stage. The most effective place to invest is the
candidate generator (for example, reading the library code that is called back,
or the package documentation), not the prompt.

\textbf{Negative results.} (i) The symbolic validator was almost neutral: the
model rarely invents package names once it sees an outline, and the alias
expansion of V2 even harmed outside names until V2.1 fixed it. (ii) The
symbolic floor was slightly harmful on the micro-benchmark: it changed 13
graphs and made 6 of them inexact; in 5 of the 13 it re-introduced calls in
code that cannot run (a constant-false branch, code after \texttt{return}, an
exception handler that cannot be reached) that the model had correctly left
out, which shows that the model reasons about control flow that the
flow-insensitive front-end ignores. On funcy the floor also re-added the
front-end's non-canonical name \code{builtins.all} for \texttt{from builtins
import all as \_all} where the model had answered the correct
\code{<builtin>.all}: V2.1's renaming of \code{builtins.X} applies to the model's
answers but not to the floor, a bug we found only while writing this report and
left in the reported results. (iii) Library hooks only paid off with
the stronger model (Fabric recall \R{fabric-in-r-nscg-v21} with Haiku,
\R{fabric-in-r-nscg-v21@sonnet} with Sonnet). (iv) A stronger model alone is not
better: Sonnet as LLM-only scores \R{macro-ins-f1-llm-only@sonnet} on the
packages, below Haiku's V0, mostly through lower recall.

\textbf{Threats to validity.}
\label{sec:threats}
\emph{Construct}: the benchmark's semantics (e.g.\ documented roles, the callback
rule, access-path names of outside values) is specific; a tool tuned to it
looks better. PyCG-2026 was developed against these benchmarks by its authors,
and our skill text restates their conventions, which favours both. Outside names
are only meaningful relative to the benchmark's naming rules, which is why the
headline is inside edges. \emph{Internal}: LLM sampling is not deterministic; we
report three runs for the main configuration, but single runs for the other
variants, so differences below about 0.02 $F_1$ are within noise. V2.1 was
designed after inspecting V2's errors on the same four packages, so its gains
over V2 are optimistic; its micro results, where it changed nothing, confirm
that the changes are small. Our scripts could contain bugs; all scoring is done by
the official scorer, and the reproduction of PyCG-2026's published numbers checks
the pipeline. \emph{External}: four packages and one benchmark family are a
small sample of Python; the packages are public, so memorisation cannot be ruled
out for the macro-benchmark (the renaming probe covers only the micro cases), and
results for other models may differ.
